تبدیلی خط کے ڈھلوان کے برابر ہے (نا کہ  ۔۔  تبدیلی خط کے ڈھلوان کی برابر ہے)
کی قیمت ڈھلوان کے برابر ہے (نا کہ  کی قیمت ڈھلوان کی برابر ہے)
مٹی کی قیمت کے  برابر ہے۔

ان جملوں میں  "کے" آزاد ہے جو مذکر اور مونث دونوں کے لئے درست ہے۔
قریب is postposition and comes with کے

کی قیمت ڈھلوان کے قریب ہو گی۔ (نا کہ   ۔۔ڈھلوان کی قریب۔۔)
کی قیمت الف کے جتنا قریب چاہیں کر سکتے ہیں (نا کہ        --کی قیمت الف کی جتنی قریب چاہیں ۔۔)

Rules:
when mentioning a variable always use masculine. eg
پانی کی گہرائی کو \عددی{h} ظاہر کرتا ہے۔
